\hypertarget{1-introduction}{%
\section{1. Introduction}\label{1-introduction}}

Supply chain resilience (SCRES) has become a central concern in
operations management as global disruptions-\/-\/-from pandemic-driven
shutdowns to geopolitical trade barriers-\/-\/-expose the fragility of
interconnected logistics networks \cite{ponomarov2009understanding,wieland2021two,hosseini2019review}. Over the past decade, reinforcement
learning (RL) has emerged as a promising paradigm for adaptive
decision-making in supply chain management (SCM). Comprehensive reviews
by Yan et al.~\cite{yan2022} and Rolf et al.~\cite{rolf2023} document a growing body of
work applying RL to inventory ordering, transportation scheduling, and
multi-echelon coordination, while Boute et al.~\cite{boute2022} provide a design
roadmap for deep RL in inventory control. More recent DES+ML reviews,
including Kogler and Maxera~\cite{kogler2026}, document papers combining discrete-event simulation (DES) with
machine learning, confirming that DES+RL is now the dominant hybrid
paradigm for operational supply chain analysis. This paper therefore
does not claim novelty in the idea of applying RL to supply chains;
instead, it contributes a benchmark-centered investigation of
\emph{when} RL works and \emph{when} it fails in resilience-oriented
operational control.

Despite this maturation, most published RL-for-SCM work concentrates on
inventory cost optimization in stylized or generic environments
\cite{gijsbrechts2022,demoor2022,dehaybe2024,geevers2024}, while the few papers addressing resilience operate
at the level of network reconfiguration rather than operational control.
Ding et al.~\cite{ding2026} represent the closest resilience competitor: they use
multi-agent PPO (MAPPO) for dynamic reconfiguration of a supply chain
system-of-systems under disruption risks, with macro-level actions such
as filling, repairing, and recruiting disrupted nodes. However, their
problem is strategic topology redesign on an abstract network,
fundamentally different from the weekly operational control-\/-\/-over
assembly shifts, inventory quantities, and transport dispatch-\/-\/-that
characterizes our setting. The gap addressed here is not whether RL can
be applied to supply chains, but whether RL can deliver resilient
operational control inside a thesis-grounded military supply chain whose
structural bottlenecks are explicitly modeled.

This work builds directly on the research agenda of Garrido-Rios~\cite{garrido2017}
and Garrido et al.~\cite{garrido2024ijpr,garrido2024iccl}. The PhD thesis developed a
13-operation military food supply chain (MFSC) simulation in
MATLAB/Simulink, establishing quantitative resilience metrics
(Re\$\^{}T\$) through 90 static configurations and data-mining
association rules. Garrido et al.~\cite{garrido2024ijpr} subsequently proposed a
Cobb-\/-Douglas factory resilience index linking production strategies
to quantitative resilience measurement. Garrido et al.~\cite{garrido2024iccl}
explicitly called for the integration of AI algorithms-\/-\/-identifying
neural networks, KANs, and reinforcement learning as the most suitable
candidates-\/-\/-into simulation-based SCRES models, arguing that
traditional DES frameworks lack endogenous learning mechanisms capable
of adaptation across successive disruption events. The present study
operationalizes this proposal: we rebuild the thesis DES in
Python/SimPy, validate it against the original reference targets
(\$-4.43\%\$ relative gap), wrap it in a Gymnasium-compatible interface,
and apply PPO-based RL control.

Our core empirical finding is dual. Under Track A---the
thesis-faithful control contract where the agent controls upstream
inventory and assembly shifts---dense common-random-number static
frontiers and multiple RL configurations across reward modes,
observation versions, and algorithms (including PPO and RecurrentPPO)
leave no publishable dynamic-control advantage: no tested learner
converts the small measured headroom over the best static policy. Under
Track B---an operational extension that adds downstream transport
dispatch control at Op10 and Op12 (an 8-dimensional action contract)
while preserving the DES backbone---PPO improves Garrido/Excel
resilience (ReT) by 0.000426 (95\% CI entirely positive) over the best
of a dense grid of 147 static dispatch comparators under paired
common-random-number evaluation, with matching gains in service
continuity and recovery-tail metrics. A 5-seed action-space ablation
shows the gain concentrates in downstream-dispatch access specifically
rather than in the number of controllable dimensions; the advantage
holds across current and increased Garrido-native risk levels at both
evaluated horizons, is mixed under severe risk; and
two independent tests rule out that the effect depends on the learned
policy exploiting privileged simulator observation fields unavailable
to static comparators. The paper's central argument is therefore
structural: RL success in this benchmark depends not on algorithmic
sophistication or reward engineering, but on whether the controllable
action space reaches the binding operational bottleneck.

The paper makes four contributions. First, we report a negative
empirical benchmark result: under the thesis-faithful action
contract, dense-frontier evaluation finds no publishable dynamic-control
advantage for learned policies. Second, we report a positive empirical
benchmark result: an operational action-space extension produces a
robust, adversarially-tested RL advantage under paired dense-frontier
evaluation. Third, we provide a structural diagnosis supported by a
causal ablation and two privileged-observation tests: the downstream
control dimensions, not action-space size or access to simulator
internals, explain the observed improvement. Fourth, we contribute a
reproducible benchmark with auditable artifacts, frozen contracts, and a
Garrido-compatible resilience evaluation pipeline. We do not claim to be
the first paper on RL for supply chains, the first DES+RL study, or to
offer architectural novelty; our contribution is experimental and
diagnostic. The remainder of the paper is organized as follows: Section
2 reviews related work across four blocks; Section 3 presents the
methodology; Section 4 reports results; Section 5 discusses implications
and limitations; and Section 6 concludes.
